\subsection*{一步一步建立直观}

\paragraph{1. 把一个猜测变成可检查的扰动}

如果猜测某个输入特征在控制输出，应固定其余输入，只改变这一项。右图的单个单元提供最简单的干预：在激活区域内，输出对第一输入的变化率就是第一权重。这个数值解释了局部敏感度，却还没有证明它在真实数据中代表某种语义。



\paragraph{2. 同一权重在不同输入处作用不同}

把输入移到关闭区域。权重没有改变，输出对小扰动却暂时不敏感，因为激活函数将这一段截平。由此可见，只阅读权重大小不足以解释网络行为；还需要输入位置、激活状态，以及后续层如何使用这个单元的输出。



\paragraph{3. 解释也必须经受反例检验}

调节偏置与目标，寻找能让原先解释失效的输入。一个“看起来检测到某特征”的单元，可能依赖背景或与该特征相关的另一信号。本窗口展示干预的逻辑，不承担真实模型的语义发现；更完整的分析还应比较样本、消融与输出变化。
